\section{问题2：资源消耗与付费行为关联分析}

\textbf{数据窗口说明：} 问题2的研究目标是资源消耗、成长速度与付费行为的关联分析，题目未限制时间窗口。本节使用30天全周期数据进行关联分析，与问题1的预测任务（仅前3天）和问题3的策略模拟（综合前两问结果）形成互补。

\subsection{资源消耗与等级增长的量化关系}

\subsubsection{分析方法}

为量化玩家每日资源消耗量与等级提升速度的关系，首先计算各玩家的日均资源消耗量：

\begin{equation}
\text{日均消耗}_r = \frac{\text{资源}r\text{的总消耗量}}{\text{生命周期天数}}
\end{equation}

其中$r \in \{\text{粮食},\text{木材},\text{矿石},\text{金币},\text{钻石}\}$。采用Pearson相关系数衡量各资源消耗与等级增长（level\_growth）之间的线性关联强度：

\begin{equation}
r_{XY} = \frac{\sum_{i=1}^{n}(X_i - \bar{X})(Y_i - \bar{Y})}{\sqrt{\sum(X_i - \bar{X})^2}\sqrt{\sum(Y_i - \bar{Y})^2}}
\label{eq:pearson}
\end{equation}

\subsubsection{结果分析}

各资源与等级增长的Pearson相关系数如下（所有$p < 0.001$）：

\begin{center}
\begin{tabular}{c|c|c|c|c|c|c}
\toprule
\textbf{资源} & 粮食 & 木材 & \textbf{矿石} & 钻石（存量） & 钻石（消耗） & 金币 \\
\midrule
$r$ & 0.394 & 0.395 & \textbf{0.504} & 0.277 & 0.379 & 0.342 \\
\bottomrule
\end{tabular}
\end{center}

矿石消耗量与等级增长的相关性最强（$r=0.504$），表明矿石是该游戏中最关键的基础资源。金币与等级增长的相关性为$r=0.342$，弱于基础资源（粮木矿），这与金币在游戏内的定位一致——金币可购买缺失资源或加速升级，但并非每日必需消耗品。钻石存量与等级增长的相关性最弱（$r=0.277$），符合钻石作为稀缺加速货币的定位——其作用不在于日常消耗，而在于突破瓶颈时刻。

图\ref{fig:p2_res}展示了各等级玩家的日均资源消耗量与日均等级增长速度（注意：此处使用日均增长率$=$等级总增长/生命周期天数，而非总增长量，以避免存活时间差异导致的伪相关）。重点发现：

\begin{enumerate}
    \item \textbf{升级瓶颈明确存在。} 1--5级玩家的日均升级速度为1.37级/天，而10级以上的玩家仅为0.37级/天，降幅达73\%。等级$>$8后日均增长速度降至0.7级/天以下，表明中后期存在显著的免费资源产出\textbf{卡点}。
    \item \textbf{资源消耗量随等级指数增长。} 10级以上玩家的日均粮食、木材、矿石消耗量分别是5级玩家的18--34倍，但等级增长仅为后者的37\%——大量资源投入未能等比例转化为升级进度。
\end{enumerate}

\begin{figure}[H]
\centering
\includegraphics[width=\textwidth]{fig5_resource_level}
\caption{各等级玩家日均资源消耗量（左）与等级增长速度（右）}
\label{fig:p2_res}
\end{figure}

\subsection{等级停滞的瓶颈分析}

\subsubsection{卡点识别}

为识别升级瓶颈的精确位置，定义\emph{等级停滞玩家}为日均升级速度$<0.1$级/天且存活$\geq 7$天的用户。按等级分段控制比较，发现瓶颈呈现两阶段特征：

\begin{table}[H]
\centering
\caption{等级瓶颈的阶段特征（停滞 vs 非停滞玩家）}
\label{tab:stagnation}
\begin{tabular}{c|c|c|c|c|c}
\toprule
\textbf{等级段} & \textbf{钻石存量比} & \textbf{矿石消耗比} & \textbf{金币存量比} & \textbf{活跃天数比} & \textbf{主因} \\
\midrule
1--3级 & 停滞:非停滞 = 1:2 & 均为0 & 均为0 & 3:1 & 钻石稀缺 \\
4--6级 & 持平(600 vs 594) & 停滞者8.2$\times$ & 停滞者$>$非停滞 & 2:1 & 资源效率递减 \\
\bottomrule
\end{tabular}
\end{table}

\textbf{早期瓶颈（1--3级）：钻石稀缺。} 停滞玩家的钻石存量仅为非停滞者的50\%（100 vs 200），其余资源消耗模式几乎相同。停滞玩家虽然活跃天数更多（3天 vs 1天），但因缺少钻石购买加速道具，升级进度无法突破。此阶段关键词：\emph{缺钻卡级}。

\textbf{中期瓶颈（4--6级）：资源使用效率极低。} 停滞玩家消耗了8.2倍矿石、4.3倍粮食/木材，活跃天数为2倍（10天 vs 5天），但日均升级速度仍低于0.1级/天。大量基础资源投入未能等比例转化为升级进度——资源边际效率随等级急剧递减。此阶段关键词：\emph{刷不动了}。

\subsubsection{金币与钻石对等级停滞的差异化角色}

为量化金币和钻石对等级停滞的因果影响，控制当前等级后建立Logistic回归：

\begin{equation}
P(\text{停滞} \mid X) = \frac{1}{1 + e^{-(\beta_0 + \beta_1\cdot\log(1+\text{钻石}) + \beta_2\cdot\log(1+\text{金币}) + \beta_3\cdot\text{等级})}}
\end{equation}

以日均增长$<0.1$级/天且存活$\geq 7$天作为停滞标签，结果：

\begin{center}
\fbox{$\beta_{\text{钻石}} = -0.355$\quad（保护因子：钻石越多，停滞概率越低）}\\[4pt]
\fbox{$\beta_{\text{金币}} = +0.157$\quad（风险因子：金币越多，停滞概率越高）}
\end{center}

\textbf{金币与钻石扮演截然不同的角色。} 钻石是\emph{瓶颈打破者}——存量越高，停滞概率越低；金币是\emph{停滞堆积物}——玩家卡级期间金币无法有效转化为升级进度，反而大量囤积。在7--9级段，停滞玩家的金币中位数高达109,898，而非停滞玩家为0——进一步验证了金币是"停滞后囤积的产物"而非"阻止停滞的手段"。

将金币存量纳入等级4+全部玩家的偏相关分析进一步确认了这一判断：金币与日均等级增长的相关性（$r=-0.143$）在五类资源中最低，远弱于矿石（$r=-0.462$）和粮食/木材（$r\approx-0.427$），说明金币无论是作为升级推动力还是瓶颈指标，解释力均十分有限。

\subsection{钻石存量与流失风险}

为研究钻石存量对流失风险的预测作用，将``7天内流失''定义为二分类目标变量（churned = 1 if lifecycle\_days $<$ 7），以对数化钻石存量中位数作为预测变量，建立Logistic回归模型：

\begin{equation}
P(\text{流失} \mid \text{钻石}) = \frac{1}{1 + e^{-(\beta_0 + \beta_1 \cdot \log(1+\text{钻石}))}}
\end{equation}

模型拟合后，求解$P(\text{流失} \mid \text{钻石}) = 0.5$时的钻石存量值作为风险阈值。结果如下：

\begin{center}
\fbox{\textbf{钻石存量$<$ 566时，7天内流失概率超过50\%}}
\end{center}

图\ref{fig:p2_diamond}直观展示了钻石存量分布与流失概率曲线。这一阈值具有直接的操作意义：当监测到玩家钻石存量跌破566时，应触发低价钻石补给礼包推送，以降低流失风险。

\begin{figure}[H]
\centering
\includegraphics[width=\textwidth]{fig6_diamond_churn}
\caption{钻石存量分布：留存vs流失玩家（左）与流失概率曲线（右）}
\label{fig:p2_diamond}
\end{figure}

\subsection{首次付费分析}

对75名付费用户（占总体的3.8\%）进行分析，主要发现：

\begin{itemize}
    \item 首付等级：100\%的付费用户首次付费发生在1级（注册初期）
    \item 付费金额分布：中位数0.99元，75\%分位数1.98元——绝大多数为6元首充档
    \item 付费金额与留存的关系：高付费组（\$100+）平均生命周期71.0天，低付费组（$<\$6$）平均26.8天\\
    付费金额越高，生命周期越长，呈现明显的正相关趋势
\end{itemize}

图\ref{fig:p2_pay}展示了付费金额分布与不同付费水平下的留存时长。

\begin{figure}[H]
\centering
\includegraphics[width=\textwidth]{fig8_first_pay}
\caption{付费金额分布（左）与不同付费水平下的平均生命周期（右）}
\label{fig:p2_pay}
\end{figure}

\subsection{XGBoost回归分析总付费关键因子}

\subsubsection{模型建立}

为分析影响总付费金额的关键因子，采用XGBoost回归模型。由于总付费金额呈极度右偏分布（96.2\%用户付费为0），模型的预测精度有限（MAE=0.45，付费者$R^2=-0.61$），因此XGBoost主要用于特征重要性解释（通过SHAP值），而非精确预测付费金额。为缓解偏态分布，对目标变量进行对数变换：

\begin{equation}
y_{\log} = \ln(1 + \text{total\_pay})
\end{equation}

选取25个特征，包括活跃度特征、等级成长特征、资源获取与消耗特征、社交特征和付费特征。训练参数：$n_{\text{estimators}}=200$，$\text{max\_depth}=5$，$\text{learning\_rate}=0.05$。

\subsubsection{特征重要性（SHAP值）}

为解释模型预测，计算SHAP（SHapley Additive exPlanations）值量化各特征对预测结果的边际贡献。图\ref{fig:p2_shap}展示了排名前12的特征（按平均|SHAP值|排序）。关键发现：

\begin{enumerate}
    \item \textbf{diamond\_flow（钻石净流量）}是最重要的付费预测因子——钻石的获取与消耗差额直接反映付费后的资源转化行为，是模型中最强的付费信号
    \item \textbf{n\_event\_types（事件类型数）}紧随其后——参与游戏活动种类越多的玩家，接触付费场景的概率越高，付费潜力越大
    \item \textbf{food\_get（粮食获取量）}和\textbf{total\_get（总资源获取量）}分列第三、四位——资源获取总量是玩家投入度的综合代理，高投入玩家更有动力为加速进度而付费
    \item \textbf{level\_end（最终等级）}位列第五——等级并非直接驱动付费，而是长期投入的自然结果；SHAP分析表明其预测贡献低于资源流动特征
\end{enumerate}

\begin{figure}[H]
\centering
\includegraphics[width=0.8\textwidth]{fig9_xgb_importance}
\caption{XGBoost特征重要性：与总付费金额关联最强的预测因子}
\label{fig:p2_shap}
\end{figure}｜｜DSML｜｜parameter>
